\documentclass{article}
\usepackage[T2A]{fontenc}
\usepackage[utf8]{inputenc} 
\usepackage[russian]{babel}
\usepackage{amsmath, amssymb, amsthm}
\usepackage{geometry}
\geometry{a4paper, margin=2cm}

\title{Минимум. Вы просили - мы сделали.}
\date{}

\begin{document}
\maketitle

\textbf{1. !!! Табличные производные и интегралы!}

\textit{Основные производные:}
\begin{align*}
(x^n)' &= n x^{n-1}, \quad (a^x)' = a^x \ln a, \quad (e^x)' = e^x, \quad (\ln x)' = \frac{1}{x} \\
(\sin x)' &= \cos x, \quad (\cos x)' = -\sin x, \quad (\tan x)' = \frac{1}{\cos^2 x}, \quad (\cot x)' = -\frac{1}{\sin^2 x} \\
(\arcsin x)' &= \frac{1}{\sqrt{1-x^2}}, \quad (\arctan x)' = \frac{1}{1+x^2}
\end{align*}

\textit{Основные интегралы:}
\begin{align*}
\int x^n dx &= \frac{x^{n+1}}{n+1} + C \ (n \neq -1), \quad \int \frac{1}{x} dx = \ln|x| + C, \quad \int e^x dx = e^x + C \\
\int \sin x dx &= -\cos x + C, \quad \int \cos x dx = \sin x + C, \quad \int \frac{dx}{x^2+a^2} = \frac{1}{a} \arctan \frac{x}{a} + C \\
\int \frac{dx}{\sqrt{a^2-x^2}} = \arcsin \frac{x}{a} + C, &\quad \int \frac{dx}{x^2-a^2} = \frac{1}{2a} \ln\left|\frac{x-a}{x+a}\right| + C, \quad \int \frac{dx}{\sqrt{x^2+q}} = \ln\left|x + \sqrt{x^2+q}\right| + C
\end{align*}

---

\textbf{2. !!! Дайте строгое определение первообразной для функции $f(x)$ на интервале (а. ).}

\textit{Ответ:} Функция $F(x)$ называется первообразной для функции $f(x)$ на интервале $(a, b)$, если $F(x)$ дифференцируема на $(a, b)$ и для всех $x \in (a, b)$ выполняется равенство $F'(x) = f(x)$.

---

\textbf{3. !!! Сформулируйте теорему о замене переменной (подстановке) в неопределенном интеграле (условия на функции и итоговую формулу).}

\textit{Ответ:} Пусть функция $x = \varphi(t)$ определена и дифференцируема на промежутке $T$, а функция $f(x)$ имеет первообразную на множестве $X = \varphi(T)$. Тогда на промежутке $T$ справедлива формула подстановки:
$$ \int f(x) dx = \int f(\varphi(t)) \varphi'(t) dt $$
где после интегрирования правой части выполняется обратная подстановка $t = \varphi^{-1}(x)$.

---

\textbf{4. !!! Запишите выражения для производной от неопределенного интеграла и для интеграла от дифференциала функции.}

\textit{Ответ:}
1) $\frac{d}{dx} \left( \int f(x) dx \right) = f(x)$ \\
2) $\int d(f(x)) = f(x) + C$

---

\textbf{5. !!! Дайте определение неопределенного интеграла. Опишите геометрическую интерпретацию семейства первообразных.}

\textit{Ответ:} Неопределенным интегралом от функции $f(x)$ называется совокупность всех ее первообразных на данном промежутке, обозначаемая $\int f(x) dx = F(x) + C$, где $F'(x) = f(x)$, а $C \in \mathbb{R}$ — произвольная постоянная. Геометрически семейство первообразных представляет собой совокупность кривых $y = F(x) + C$, которые получаются друг из друга параллельным переносом вдоль оси $Oy$.

---

\textbf{6. !!! Сформулируйте теорему об общем виде первообразных для функции на интервале.}

\textit{Ответ:} Если $F(x)$ является одной из первообразных для функции $f(x)$ на интервале $(a, b)$, то любая другая первообразная $\Phi(x)$ этой функции на данном интервале имеет вид $\Phi(x) = F(x) + C$, где $C$ — некоторая вещественная постоянная.

---

\textbf{7. !!! Сформулируйте свойство линейности неопределенного интеграла.}

\textit{Ответ:} Если функции $f(x)$ и $g(x)$ имеют первообразные на некотором промежутке, а $\alpha, \beta \in \mathbb{R}$ — произвольные числа ($\alpha^2 + \beta^2 \neq 0$), то функция $\alpha f(x) + \beta g(x)$ также имеет первообразную на этом промежутке, и справедливо равенство:
$$ \int (\alpha f(x) + \beta g(x)) dx = \alpha \int f(x) dx + \beta \int g(x) dx $$

---

\textbf{8. !!! Сформулируйте теорему о разложении правильной рациональной дроби на простейшие (виды слагаемых для кратных линейных и квадратичных множителей).}

\textit{Ответ:} Любую правильную рациональную дробь $\frac{P(x)}{Q(x)}$ со знаменателем, разложенным на линейные и неприводимые квадратичные множители вида $Q(x) = (x-a)^k \dots (x^2+px+q)^m$ (где $p^2-4q < 0$), можно единственным образом представить в виде суммы простейших дробей. При этом каждому множителю $(x-a)^k$ соответствует группа из $k$ слагаемых вида:
$$ \frac{A_1}{x-a} + \frac{A_2}{(x-a)^2} + \dots + \frac{A_k}{(x-a)^k} $$
а каждому квадратичному множителю $(x^2+px+q)^m$ соответствует группа из $m$ слагаемых вида:
$$ \frac{M_1 x + N_1}{x^2+px+q} + \frac{M_2 x + N_2}{(x^2+px+q)^2} + \dots + \frac{M_m x + N_m}{(x^2+px+q)^m} $$

---

\textbf{9. !!! Запишите канонический вид элементарных (простейших) дробей І, ІІ, ІІІ и IV типов.}

\textit{Ответ:}
\begin{align*}
\text{I тип:} & \quad \frac{A}{x-a} \\
\text{II тип:} & \quad \frac{A}{(x-a)^k}, \quad (k \ge 2, k \in \mathbb{N}) \\
\text{III тип:} & \quad \frac{Mx+N}{x^2+px+q}, \quad (p^2-4q < 0) \\
\text{IV тип:} & \quad \frac{Mx+N}{(x^2+px+q)^m}, \quad (m \ge 2, m \in \mathbb{N}, p^2-4q < 0)
\end{align*}

---

\textbf{10. Укажите общий вид интеграла от элементарной дроби III типа и выпишите формулу выделения полного квадрата в трехчлене $x^{2}+px+q$.}

\textit{Ответ:} Общий вид интеграла: $\int \frac{Mx+N}{x^2+px+q} dx$. Формула выделения полного квадрата:
$$ x^2+px+q = \left(x + \frac{p}{2}\right)^2 + \left(q - \frac{p^2}{4}\right) $$

---

\textbf{11. Запишите формулу универсальной тригонометрической подстановки, а также выражения для sinx, cosx и через новый параметр $t=tg(x/2)$.}

\textit{Ответ:} Подстановка: $t = \tan \frac{x}{2}$. Выражения для тригонометрических функций и дифференциала:
$$ \sin x = \frac{2t}{1+t^2}, \quad \cos x = \frac{1-t^2}{1+t^2}, \quad dx = \frac{2dt}{1+t^2} $$

---

\textbf{12. Запишите общее алгебраическое уравнение-шаблон метода Остроградского, четко указав структуры знаменателей рациональной и трансцендентной частей.}

\textit{Ответ:} Уравнение-шаблон:
$$ \int \frac{P(x)}{Q(x)} dx = \frac{X(x)}{Q_1(x)} + \int \frac{Y(x)}{Q_2(x)} dx $$
где $Q_2(x) = \text{НОД}(Q(x), Q'(x))$ — многочлен, содержащий все корни многочлена $Q(x)$, взятые по одному разу; $Q_1(x) = \frac{Q(x)}{Q_2(x)}$ — многочлен, содержащий кратные корни $Q(x)$ со степенями, уменьшенными на единицу; $X(x)$ и $Y(x)$ — многочлены с неопределенными коэффициентами, степени которых на единицу меньше степеней многочленов $Q_1(x)$ и $Q_2(x)$ соответственно.

---

\textbf{13. Перечислите случаи, когда применение метода Остроградского является единствен- но целесообразным.}

\textit{Ответ:} Применение метода Остроградского целесообразно тогда, когда знаменатель $Q(x)$ правильной рациональной дроби имеет вещественные или комплексные корни высокой кратности, прямое интегрирование которых через стандартное разложение на простейшие дроби потребовало бы громоздкого вычисления интегралов IV типа методом рекуррентных соотношений.

---

\textbf{14. Укажите общий вид и алгоритм дробно-линейной подстановки для рационализации интегралов, содержащих радикалы вида $\sqrt[n]{\frac{ax+b}{cx+d}}$ где ad $-bc\ne0$.}

\textit{Ответ:} Общий вид интеграла: $\int R\left(x, \sqrt[n]{\frac{ax+b}{cx+d}}\right) dx$. \\
Алгоритм подстановки:
1) Вводится новая переменная $t = \sqrt[n]{\frac{ax+b}{cx+d}}$, откуда $t^n = \frac{ax+b}{cx+d}$. \\
2) Выражается переменная $x$: $x = \frac{dt^n - b}{a - ct^n}$. \\
3) Находится дифференциал $dx = \left(\frac{dt^n - b}{a - ct^n}\right)' dt$. \\
4) Полученные выражения подставляются под знак интеграла, что сводит его к интегралу от рациональной функции переменной $t$.

---

\textbf{15. Запишите формулу замены переменной и выражения для х и , если под знаком интеграла присутствует радикал $\sqrt{x^{2}+q}$.}

\textit{Ответ:} Для избавления от радикала используется гиперболическая подстановка (при $q > 0$):
$$ x = \sqrt{q} \sinh t \implies dx = \sqrt{q} \cosh t \, dt, \quad \sqrt{x^2+q} = \sqrt{q} \cosh t $$
Или тригонометрическая подстановка:
$$ x = \sqrt{q} \tan t \implies dx = \frac{\sqrt{q}}{\cos^2 t} dt, \quad \sqrt{x^2+q} = \frac{\sqrt{q}}{\cos t} $$

---

\textbf{16. Запишите формулировки всех трех классических подстановок Эйлера для выражения $\sqrt{ax^{2}+bx+c}$, указав алгебраические условия применимости для каждой из них.}

\textit{Ответ:}
1) \textbf{Первая подстановка Эйлера:} применяется при $a > 0$:
$$ \sqrt{ax^2+bx+c} = \pm\sqrt{a}x + t $$
2) \textbf{Вторая подстановка Эйлера:} применяется при $c > 0$:
$$ \sqrt{ax^2+bx+c} = xt \pm \sqrt{c} $$
3) \textbf{Третья подстановка Эйлера:} применяется, если дискриминант $D = b^2 - 4ac > 0$, то есть квадратный трехчлен имеет вещественные корни $x_1$ и $x_2$:
$$ \sqrt{ax^2+bx+c} = t(x - x_1) $$

---

\textbf{17. Запишите стандартный вид дифференциального бинома Чебышёва и укажите, какими множествами должны быть его параметры а, , т, п.р.}

\textit{Ответ:} Стандартный вид: $x^m (a + b x^n)^p dx$. Параметры: $a, b \in \mathbb{R} \setminus \{0\}$, а показатели степеней $m, n, p \in \mathbb{Q}$ (множество рациональных чисел).

---

\textbf{18. Сформулируйте три точных условия теоремы Чебышёва, при которых интеграл от дифференциального бинома выражается в элементарных функциях.}

\textit{Ответ:} Интеграл $\int x^m (a + b x^n)^p dx$ выражается в элементарных функциях только в следующих трех случаях:
1) $p \in \mathbb{Z}$ (целое число). \\
2) $\frac{m+1}{n} \in \mathbb{Z}$ (целое число). \\
3) $\frac{m+1}{n} + p \in \mathbb{Z}$ (целое число).

---

\textbf{19. Укажите конкретные подстановки, соответствующие каждому из трех случаев интегрируемости по Чебышёву.}

\textit{Ответ:}
1) Если $p \in \mathbb{Z}$, подстановка: $x = t^s$, где $s$ — общий знаменатель дробей $m$ и $n$. \\
2) Если $\frac{m+1}{n} \in \mathbb{Z}$, подстановка: $a + b x^n = t^s$, где $s$ — знаменатель дроби $p$. \\
3) Если $\frac{m+1}{n} + p \in \mathbb{Z}$, подстановка: $a x^{-n} + b = t^s$ (или $a + bx^n = t^s x^n$), где $s$ — знаменатель дроби $p$.

---

\textbf{20. !!! Запишите определение интегральной суммы Римана $\sigma(f,T,\{\xi_{i}\})$ для функции $f(x)$ на сегменте [ , b], указав структуру разбиения Т и выборки точек $\{\xi_{i}\}$.}

\textit{Ответ:} Пусть на сегменте $[a, b]$ задано разбиение $T: a = x_0 < x_1 < \dots < x_n = b$, разделяющее его на частичные отрезки длины $\Delta x_i = x_i - x_{i-1}$. На каждом отрезке $[x_{i-1}, x_i]$ выбрана точка $\xi_i \in [x_{i-1}, x_i]$. Интегральной суммой Римана называется величина:
$$ \sigma(f, T, \{\xi_i\}) = \sum_{i=1}^n f(\xi_i) \Delta x_i $$

---

\textbf{21. !!! Сформулируйте на языке $\epsilon-\delta$ понятие предела интегральных сумм при мелкости разбиения $\lambda(T)\rightarrow0$.}

\textit{Ответ:} Число $I$ называется пределом интегральных сумм Римана при мелкости разбиения $\lambda(T) = \max_{1 \le i \le n} \Delta x_i \to 0$, если:
$$ \forall \epsilon > 0 \quad \exists \delta > 0 \quad \forall T \, (\lambda(T) < \delta) \quad \forall \{\xi_i\} \implies |\sigma(f, T, \{\xi_i\}) - I| < \epsilon $$

---

\textbf{22. !!! Сформулируйте необходимое условие интегрируемости функции по Риману на сегменте и приведите контрпример, доказывающий, что оно не является достаточ- ным.}

\textit{Ответ:} Необходимое условие: функция $f(x)$ должна быть ограниченной на сегменте $[a, b]$. \\
Контрпример: Функция Дирихле $D(x) = \begin{cases} 1, & x \in \mathbb{Q} \\ 0, & x \notin \mathbb{Q} \end{cases}$ на сегменте $[0, 1]$. Она ограничена ($|D(x)| \le 1$), но не является интегрируемой по Риману, так как любая верхняя сумма Дарбу равна 1, а нижняя равна 0.

---

\textbf{23. !!! Запишите определения нижней $S_{*}(T)$ и верхней $S^{*}(T)$ сумм Дарбу, используя точные грани $m_{i}$ и $M_{i}$ функции на частичных отрезках.}

\textit{Ответ:} Пусть $m_i = \inf_{x \in [x_{i-1}, x_i]} f(x)$ и $M_i = \sup_{x \in [x_{i-1}, x_i]} f(x)$. Тогда:
$$ \text{Нижняя сумма Дарбу:} \quad S_*(T) = \sum_{i=1}^n m_i \Delta x_i $$
$$ \text{Верхняя сумма Дарбу:} \quad S^*(T) = \sum_{i=1}^n M_i \Delta x_i $$

---

\textbf{24. !!! Дайте определение нижнего (1) и верхнего $(\overline{I})$ интегралов Дарбу. Какое неравенство всегда связывает эти две величины?}

\textit{Ответ:} Нижний интеграл Дарбу $I_*$ — это точная верхняя грань нижних сумм Дарбу по всем возможным разбиениям $T$: $I_* = \sup_{T} S_*(T)$. Верхний интеграл Дарбу $\overline{I}$ — это точная нижняя грань верхних сумм Дарбу по всем возможным разбиениям $T$: $\overline{I} = \inf_{T} S^*(T)$. Между ними всегда выполняется неравенство:
$$ I_* \le \overline{I} $$

---

\textbf{25. Что такое колебание $\omega_{i}$ функции $f(x)$ на сегменте $[x_{i-1},x_{i}]$? Как оно связано с точными гранями функции?}

\textit{Ответ:} Колебанием $\omega_i$ функции $f(x)$ на отрезке $[x_{i-1}, x_i]$ называется разность между ее точной верхней и точной нижней гранями на этом отрезке:
$$ \omega_i = M_i - m_i = \sup_{x \in [x_{i-1}, x_i]} f(x) - \inf_{x \in [x_{i-1}, x_i]} f(x) $$

---

\textbf{26. !!! Сформулируйте основной критерий интегрируемости ограниченной функции в терминах совпадения интегралов Дарбу.}

\textit{Ответ:} Ограниченная на сегменте $[a, b]$ функция $f(x)$ интегрируема по Риману тогда и только тогда, когда ее нижний и верхний интегралы Дарбу равны между собой:
$$ I_* = \overline{I} $$

---

\textbf{27. Сформулируйте критерий интегрируемости Римана в терминах колебаний функции на разбиении.}

\textit{Ответ:} Ограниченная функция $f(x)$ интегрируема на сегменте $[a, b]$ тогда и только тогда, когда предел суммы произведений колебаний на длины частичных отрезков равен нулю при стремлении мелкости разбиения к нулю:
$$ \lim_{\lambda(T) \to 0} \sum_{i=1}^n \omega_i \Delta x_i = 0 $$

---

\textbf{28. Запишите геометрический смысл разности между верхней и нижней суммами Дарбу $S^{*}(T)-S_{*}(T)$ в терминах площади ступенчатых фигур.}

\textit{Ответ:} Разность $S^*(T) - S_*(T)$ геометрически равна площади ступенчатой фигуры, составленной из прямоугольников с основаниями $\Delta x_i$ и высотами $\omega_i$. Эта величина выражает площадь области между объемлющей (описанной) и вписанной ступенчатыми фигурами, покрывающими график функции.

---

\textbf{29. !!! Дайте определение равномерной непрерывности функции $f(x)$ на множестве $X$ (на языке $\epsilon$-$\delta$).}

\textit{Ответ:} Функция $f(x)$ называется равномерно непрерывной на множестве $X$, если для любого $\varepsilon > 0$ существует такое $\delta > 0$, что для любых точек $x_1, x_2 \in X$, удовлетворяющих условию $|x_1 - x_2| < \delta$, выполняется неравенство $|f(x_1) - f(x_2)| < \varepsilon$:
$$ \forall \varepsilon > 0 \quad \exists \delta > 0 \quad \forall x_1, x_2 \in X : \quad |x_1 - x_2| < \delta \implies |f(x_1) - f(x_2)| < \varepsilon $$

---
---

\textbf{30. !!! Сформулируйте великую теорему Кантора о равномерной непрерывности функ- ций на компакте.}

\textit{Ответ:} Если функция $f(x)$ непрерывна на компактном множестве $X$ (в частности, на замкнутом сегменте $[a, b]$), то она равномерно непрерывна на нем.

---

\textbf{31. !!! Запишите формулировку критерия интегрируемости Римана в терминах разности верхней и нижней сумм Дарбу $(U(T)-L(T)<\epsilon)$.}

\textit{Ответ:} Ограниченная на сегменте $[a, b]$ функция $f(x)$ интегрируема по Риману тогда и только тогда, когда для любого $\epsilon > 0$ существует такое разбиение $T$, что разность между ее верхней и нижней суммами Дарбу меньше $\epsilon$:
$$ S^*(T) - S_*(T) < \epsilon $$

---

\textbf{32. Дайте определение монотонной (неубывающей /невозрастающей) функции на сегмен- те [а, 6]. Обязана ли монотонная функция быть непрерывной?}

\textit{Ответ:} Функция $f(x)$ называется неубывающей на $[a, b]$, если $\forall x_1, x_2 \in [a, b] : x_1 < x_2 \implies f(x_1) \le f(x_2)$ (невозрастающей, если $f(x_1) \ge f(x_2)$). Монотонная функция не обязана быть непрерывной (например, функция разрывного скачка $f(x) = \lfloor x \rfloor$).

---

\textbf{33. Какое свойство ограниченности автоматически вытекает из факта монотонности функции на сегменте [ , b]? Запишите точные грани функции.}

\textit{Ответ:} Из факта монотонности функции на сегменте автоматически вытекает ее ограниченность. Если функция $f(x)$ неубывающая на $[a, b]$, ее точные грани равны значениям на границах сегмента:
$$ \inf_{x \in [a, b]} f(x) = f(a), \quad \sup_{x \in [a, b]} f(x) = f(b) $$

---

\textbf{34. Приведите пример монотонной функции, имеющей бесконечное (счетное) число то- чек разрыва на сегменте.}

\textit{Ответ:} Функция «кусочная лестница», построенная на основе перенумерованных рациональных точек $q_n$ сегмента $[a, b]$:
$$ f(x) = \sum_{q_n \le x} \frac{1}{2^n} $$
Данная функция строго монотонно возрастает и имеет разрывы в каждой рациональной точке.

---

\textbf{35. !!! Дайте определение точки разрыва первого рода и точки разрыва второго рода для вещественной функции.}

\textit{Ответ:} Точка $x_0$ называется точкой разрыва первого рода, если в ней существуют конечные односторонние пределы $f(x_0 - 0)$ и $f(x_0 + 0)$, но они не равны друг другу или не равны значению $f(x_0)$. Точка $x_0$ называется точкой разрыва второго рода, если хотя бы один из односторонних пределов $f(x_0 - 0)$ или $f(x_0 + 0)$ не существует или равен бесконечности.

---

\textbf{36. !!! Сформулируйте теорему об интегрируемости функций с конечным числом точек разрыва на сегменте.}

\textit{Ответ:} Если ограниченная на сегменте $[a, b]$ функция $f(x)$ имеет конечное число точек разрыва, то она интегрируема на этом сегменте по Риману.

---

\textbf{37. Рассмотрите функцию $f(x)=1/x$ на полуинтервале $(0,1]u~f(0)=0$ Является ли она интегрируемой по Риману на [0, 1]? Определите на основе условий интегрируемости.}

\textit{Ответ:} Нет, не является. Функция $f(x)$ неограничена в окрестности точки $x=0$ ($\lim_{x \to 0^+} \frac{1}{x} = +\infty$). Нарушение необходимого условия ограниченности исключает интегрируемость по Риману.

---

\textbf{38. !!! Сформулируйте свойство линейности определенного интеграла (интеграл линей- ной комбинации функций).}

\textit{Ответ:} Если функции $f(x)$ и $g(x)$ интегрируемы на сегменте $[a, b]$, то для любых чисел $\alpha, \beta \in \mathbb{R}$ их линейная комбинация также интегрируема на $[a, b]$, и выполняется равенство:
$$ \int_a^b (\alpha f(x) + \beta g(x)) dx = \alpha \int_a^b f(x) dx + \beta \int_a^b g(x) dx $$

---

\textbf{39. Сформулируйте свойство сохранения знака для определенного интеграла от неотри- цательной функции.}

\textit{Ответ:} Если функция $f(x)$ интегрируема на сегменте $[a, b]$ ($a < b$) и $f(x) \ge 0$ для всех $x \in [a, b]$, то определенный интеграл от этой функции также неотрицателен:
$$ \int_a^b f(x) dx \ge 0 $$

---

\textbf{40. !!! Сформулируйте свойство аддитивности определенного интеграла по промежутку интегрирования для случая $a<c<b$.}

\textit{Ответ:} Если функция $f(x)$ интегрируема на отрезках $[a, c]$ и $[c, b]$, то она интегрируема и на отрезке $[a, b]$, причем:
$$ \int_a^b f(x) dx = \int_a^c f(x) dx + \int_c^b f(x) dx $$

---

\textbf{41. !!! Распространите определение интеграла $\int_{a}^{b}f(x)$ dx на случаи, когда $a=b$ и $a>b$. Сохраняется ли при этом формальная алгебраическая аддитивность?}

\textit{Ответ:} По определению полагают $\int_a^a f(x) dx = 0$, а для случая $a > b$: $\int_a^b f(x) dx = -\int_b^a f(x) dx$. При таких определениях формальная алгебраическая аддитивность $\int_a^b f = \int_a^c f + \int_c^b f$ полностью сохраняется для любого взаимного расположения точек $a, b, c$.

---

\textbf{42. В чем заключается физический смысл аддитивности интеграла (например, в терми- нах массы неоднородного стержня или работы переменной силы)?**}

\textit{Ответ:} Физический смысл заключается в свойстве аддитивности результирующей величины: полная масса неоднородного стержня складывается из масс его отдельных частей, а суммарная работа переменной силы на траектории равна сумме работ, совершенных на последовательных участках пути.

---

\textbf{43. !!! Сформулируйте теорему об интегрировании неравенств $(ecin~f(x)\le g(x)$ на [ , b]).}

\textit{Ответ:} Если функции $f(x)$ и $g(x)$ интегрируемы на сегменте $[a, b]$ ($a < b$) и для всех $x \in [a, b]$ справедливо неравенство $f(x) \le g(x)$, то:
$$ \int_a^b f(x) dx \le \int_a^b g(x) dx $$

---

\textbf{44. !!! Запишите неравенство, связывающее модуль определенного интеграла и опреде- ленный интеграл от модуля функции.}

\textit{Ответ:} Если $f(x) \in \mathcal{R}[a, b]$ ($a < b$), то выполняется неравенство:
$$ \left| \int_a^b f(x) dx \right| \le \int_a^b |f(x)| dx $$

---

\textbf{45. Сформулируйте необходимое и достаточное условие существования интеграла от мо- дуля $|f(x)|$, если известно, что $f(x)\in\mathcal{R}[a,b]$. Верно ли обратное утверждение?}

\textit{Ответ:} Если $f(x)$ интегрируема по Риману на $[a, b]$, то интеграл от ее модуля $\int_a^b |f(x)| dx$ всегда существует (это гарантируется неравенством для колебаний $\omega_i(|f|) \le \omega_i(f)$). Обратное утверждение неверно: для разрывной модифицированной функции Дирихле ($f(x) = 1$ в рациональных точках, $f(x) = -1$ в иррациональных) интеграл не существует, однако $|f(x)| \equiv 1$ является интегрируемой функцией.

---

\textbf{46. !!! Сформулируйте первую теорему о среднем значении определенного интеграла Римана для непрерывной функции.}

\textit{Ответ:} Если функция $f(x)$ непрерывна на сегменте $[a, b]$, то существует точка $\xi \in [a, b]$ такая, что выполняется равенство:
$$ \int_a^b f(x) dx = f(\xi)(b - a) $$

---

\textbf{47. Запишите формулу для вычисления среднего значения и интегрируемой функции $f(x)$ на сегменте [а, ].}

\textit{Ответ:} Среднее значение $\mu$ функции $f(x)$ на сегменте $[a, b]$ вычисляется по формуле:
$$ \mu = \frac{1}{b - a} \int_a^b f(x) dx $$

---

\textbf{48. !!! Сформулируйте теорему Вейерштрасса о достижении непрерывной функцией сво- их точных граней на сегменте (база для теоремы о среднем).}

\textit{Ответ:} Функция $f(x)$, непрерывная на замкнутом сегменте $[a, b]$, ограничена на нем и достигает своих точных верхней и нижней граней (то есть существуют точки $x_1, x_2 \in [a, b]$ такие, что $f(x_1) = \inf f(x)$ и $f(x_2) = \sup f(x)$).

---

\textbf{49. !!! Сформулируйте обобщенную теорему о среднем значении для произведения двух функций $f(x)$ и $g(x)$.}

\textit{Ответ:} Если функция $f(x)$ непрерывна на сегменте $[a, b]$, а функция $g(x)$ интегрируема и сохраняет знак (неотрицательна или неположительна) на этом сегменте, то существует точка $\xi \in [a, b]$ такая, что:
$$ \int_a^b f(x)g(x) dx = f(\xi) \int_a^b g(x) dx $$

---

\textbf{50. Какое ограничение накладывается на весовую функцию $g(x)$ в обобщенной теореме о среднем? Сохранится ли структура теоремы без этого ограничения?}

\textit{Ответ:} На функцию $g(x)$ накладывается строгое ограничение сохранения знака на всем сегменте интегрирования. Без этого ограничения структура теоремы разрушается, так как интеграл от произведения может выдать значение, выводящее величину $f(\xi)$ за пределы области значений функции $f(x)$.

---

\textbf{51. !!! Запишите формулировку теоремы Коши (о промежуточном значении) для непре- рывной функции, принимающей любые значения между своими точными гранями ти М.}

\textit{Ответ:} Если функция $f(x)$ непрерывна на сегменте $[a, b]$, $m = \min f(x)$ и $M = \max f(x)$, то для любого числа $C \in [m, M]$ существует точка $\xi \in [a, b]$ такая, что выполняется равенство $f(\xi) = C$.

---

\textbf{52. !!! Запишите определение функции интеграла с переменным верхним пределом $\Phi(x)$. Укажите её область определения, если исходная функция $f\in\mathcal{R}[a,b]$.}

\textit{Ответ:} Функция интеграла с переменным верхним пределом определяется формулой:
$$ \Phi(x) = \int_a^x f(t) dt $$
Если $f \in \mathcal{R}[a, b]$, то областью определения функции $\Phi(x)$ является весь сегмент $[a, b]$.

---

\textbf{53. !!! Сформулируйте теорему о непрерывности интеграла с переменным верхним пре- делом.}

\textit{Ответ:} Если функция $f(x)$ интегрируема на сегменте $[a, b]$, то функция $\Phi(x) = \int_a^x f(t) dt$ непрерывна на этом сегменте (и более того, принадлежит классу функций, удовлетворяющих условию Липшица).

---

\textbf{54. Если функция $f(x)$ имеет на сегменте [ , b] одну точку разрыва первого рода, будет ли функция $\Phi(x)$ непрерывной в этой точке? Обоснуйте в одну фразу.}

\textit{Ответ:} Да, функция $\Phi(x)$ будет непрерывной в этой точке, так как интеграл от любой ограниченной функции всегда является непрерывной функцией своего верхнего предела, независимо от наличия изолированных точек разрыва у подынтегральной функции.

---

\textbf{55. !!! Сформулируйте теорему Барроу о дифференцируемости интеграла с перемен- ным верхним пределом. Обязательна ли непрерывность подынтегральной функции во всех точках сегмента?}

\textit{Ответ:} Если функция $f(x)$ интегрируема на $[a, b]$ и непрерывна в некоторой точке $x_0 \in [a, b]$, то функция $\Phi(x) = \int_a^x f(t) dt$ дифференцируема в этой точке, причем $\Phi'(x_0) = f(x_0)$. Непрерывность функции $f(x)$ во всех точках сегмента не обязательна — дифференцируемость интеграла гарантируется локально в каждой конкретной точке непрерывности подынтегральной функции.

---

\textbf{56. Запишите формулу дифференцирования интеграла, у которого и верхний, и нижний пределы являются дифференцируемыми функциями: $\frac{d}{dx}\int_{\psi(x)}^{\phi(x)}f(t)dt$.}

\textit{Ответ:} При условии непрерывности функции $f(t)$ формула имеет вид:
$$ \frac{d}{dx}\int_{\psi(x)}^{\phi(x)}f(t)dt = f(\phi(x))\phi'(x) - f(\psi(x))\psi'(x) $$

---

\textbf{57. Чему равна производная функции $\Phi(x)=\int_{0}^{x}|t|dt$ в точке $x=0$? Дайте краткое обоснование.}

\textit{Ответ:} Производная равна $0$. Обоснование: подынтегральная функция $f(t) = |t|$ является непрерывной в точке $t = 0$, следовательно, по теореме Барроу, производная интеграла по верхнему пределу равна значению подынтегральной функции в этой точке: $\Phi'(0) = |0| = 0$.

---

\textbf{58. Сформулируйте условия, при которых возможно применение правила Лопиталя к отношению двух интегралов с переменными верхними пределами при $x\rightarrow x_{0}$.}

\textit{Ответ:} Применение правила Лопиталя возможно, если подынтегральные функции непрерывны в окрестности точки $x_0$, пределы интегрирования при $x \to x_0$ порождают неопределенность вида $\left[\frac{0}{0}\right]$ или $\left[\frac{\infty}{\infty}\right]$, а производная интеграла, стоящего в знаменателе, не обращается в нуль в проколотой окрестности точки $x_0$.

---

\textbf{59. Запишите формулировку теоремы о существовании первообразной для непрерывной на сегменте функции. Как эта теорема связана с теоремой Барроу?}

\textit{Ответ:} Всякая функция, непрерывная на сегменте $[a, b]$, имеет первообразную на этом сегменте. Данная теорема является прямым следствием теоремы Барроу, поскольку сам интеграл с переменным верхним пределом $\Phi(x) = \int_a^x f(t) dt$ выступает в роли точной конструктивной первообразной для непрерывной функции $f(x)$.

---

\textbf{60. Вычислить предел $lim_{x\rightarrow0^{+}}\frac{\int_{0}^{x}sin(t^{2})dt}{x^{3}}$ устно, используя эквивалентные бесконечно малые под знаком интеграла.}

\textit{Ответ:} Так как при $t \to 0$ выполнено $\sin(t^2) \sim t^2$, то интеграл эквивалентен величине $\int_0^x t^2 dt = \frac{x^3}{3}$. Предел отношения равен:
$$ \lim_{x \to 0^+} \frac{\frac{1}{3}x^3}{x^3} = \frac{1}{3} $$

---

\textbf{61. !!! Сформулируйте теорему Ньютона-Лейбница.}

\textit{Ответ:} Если функция $f(x)$ непрерывна на сегменте $[a, b]$, а функция $F(x)$ является ее произвольной первообразной на этом сегменте, то имеет место формула:
$$ \int_a^b f(x) dx = F(b) - F(a) $$

---

\textbf{62. Опишите разницу между определенным интегралом Римана $\int_{a}^{b}f(x)$ dx и неопределенным интегралом $\int f(x)$ dx с точки зрения их математической природы (что яв- ляется числом, а что семейством функций?).}

\textit{Ответ:} Определенный интеграл Римана $\int_a^b f(x) dx$ по своей природе является конкретным действительным числом (геометрически — площадью криволинейной трапеции). Неопределенный интеграл $\int f(x) dx$ является множеством (семейством) функций вида $F(x) + C$, отличающихся друг от друга на произвольную постоянную.

---

\textbf{63. Справедлива ли формула Ньютона-Лейбница для функции $f(x)=sign(x)$ на сегмен- те [-1,1]? Поставьте краткий диагноз.}

\textit{Ответ:} Нет, не справедлива. Функция $f(x) = \text{sign}(x)$ имеет точку разрыва первого рода в $x=0$, вследствие чего у нее отсутствует первообразная на всем сегменте $[-1, 1]$ (любая первообразная должна быть непрерывной функцией, а функция $|x|$ не имеет производной в нуле).

---

\textbf{64. !!! Запишите интегральную формулу для вычисления площади плоской фигуры в полярных координатах, ограниченной лучами $\phi=\alpha.$ $\phi=\beta$ и непрерывной кривой $\rho=\rho(\phi)$.}

\textit{Ответ:} Формула площади имеет вид:
$$ S = \frac{1}{2} \int_{\alpha}^{\beta} \rho^2(\phi) d\phi $$

---

\textbf{65. Укажите условия применимости интегральной формулы площади в полярных коор- динатах (требования к знаку и гладкости функции $\rho(\phi))$.}

\textit{Ответ:} Функция $\rho(\phi)$ должна быть непрерывной на отрезке $[\alpha, \beta]$ и неотрицательной ($\rho(\phi) \ge 0$).

---

\textbf{66. Запишите формулу площади кругового сектора радиуса R с центральным углом $\Delta\phi$, которая служит базовым элементом аппроксимации.}

\textit{Ответ:} Формула площади элементарного сектора:
$$ \Delta S = \frac{1}{2} R^2 \Delta\phi $$

---

\textbf{67. !!! Дайте определение длины дуги кривой через точную верхнюю грань (или предел) длин вписанных ломаных.}

\textit{Ответ:} Длиной дуги кривой называется точная верхняя грань длин вписанных в эту дугу ломаных линий по всем возможным разбиениям дуги:
$$ L = \sup_{T} \sum_{i=1}^n |P_{i-1}P_i| $$

---

\textbf{68. Какая кривая на плоскости называется спрямляемой? Сформулируйте достаточное условие спрямляемости кривой (через гладкость).}

\textit{Ответ:} Кривая называется спрямляемой, если множество длин вписанных в нее ломаных ограничено сверху (то есть длина кривой конечна). Достаточным условием спрямляемости является гладкость кривой — непрерывная дифференцируемость функций, задающих кривую (класс $C^1$).

---

\textbf{69. !!! Запишите три канонические формулы для вычисления длины дуги кривой: в декартовых координатах $(y=y(x))$, в параметрической форме $(x=x(t)$, $y=y(t))$ и в полярной системе $(\rho=\rho(\phi))$.}

\textit{Ответ:}
1) Декартовы координаты: $L = \int_a^b \sqrt{1 + (y'(x))^2} dx$ \\
2) Параметрическая форма: $L = \int_{t_1}^{t_2} \sqrt{(x'(t))^2 + (y'(t))^2} dt$ \\
3) Полярная система: $L = \int_{\alpha}^{\beta} \sqrt{\rho^2(\phi) + (\rho'(\phi))^2} d\phi$

---

\textbf{70. !!! Запишите формулу для вычисления объема тела, образованного вращением кри- волинейной трапеции под графиком $y=f(x)\ge0$ вокруг оси Ох.}

\textit{Ответ:} Формула объема вращения вокруг оси $Ox$:
$$ V_x = \pi \int_a^b f^2(x) dx $$

---

\textbf{71. Запишите интегральную формулу вычисления объема тела при вращении той же трапеции вокруг оси Оу $(a\ge0)$.}

\textit{Ответ:} Формула объема вращения вокруг оси $Oy$ (метод цилиндрических оболочек):
$$ V_y = 2\pi \int_a^b x f(x) dx $$

---

\textbf{72. Дайте краткую геометрическую интерпретацию элементарного объема dV для мето- да дисков и метода цилиндрических оболочек.}

\textit{Ответ:} Для метода дисков элемент объема $dV = \pi f^2(x) dx$ представляет собой объем бесконечно тонкого круглого цилиндра (диска) радиуса $f(x)$ и толщины $dx$. Для метода цилиндрических оболочек элемент объема $dV = 2\pi x f(x) dx$ представляет собой объем бесконечно тонкого полого цилиндра (оболочки) радиуса $x$, высоты $f(x)$ и толщины $dx$.

---

\textbf{73. !!! Дайте определения абсолютной и условной сходимости для несобственного инте- грала первого рода $\int_{a}^{+\infty}f(x)dx$.}

\textit{Ответ:} Несобственный интеграл $\int_a^{+\infty} f(x) dx$ называется абсолютно сходящимся, если сходится несобственный интеграл от модуля подынтегральной функции: $\int_a^{+\infty} |f(x)| dx$. Интеграл называется условно сходящимся, если сам интеграл $\int_a^{+\infty} f(x) dx$ сходится, но интеграл от его модуля $\int_a^{+\infty} |f(x)| dx$ расходится.

---

\textbf{74. Сформулируйте признак Абеля сходимости несобственного интеграла (условия на функции $f(x)_{H}g(x))$.}

\textit{Ответ:} Интеграл $\int_a^{+\infty} f(x)g(x) dx$ сходится, если выполнены следующие условия:
1) Интеграл $\int_a^{+\infty} f(x) dx$ сходится. \\
2) Функция $g(x)$ монотонна и ограничена на полуинтервале $[a, +\infty)$.

---

\textbf{75. Запишите точное значение интеграла Дирихле $I=\int_{0}^{+\infty}\frac{sin~x}{x}dx$ . Сходится ли он аб- солютно?}

\textit{Ответ:} Точное значение интеграла Дирихле:
$$ I = \frac{\pi}{2} $$
Абсолютно данный интеграл не сходится (сходится только условно).

---

\textbf{76. !!! Дайте определение №-й частичной суммы $S_{N}$ и суммы числового ряда $\sum_{n=1}^{\infty}a_{n}.$ Что означает фраза «ряд сходится»?}

\textit{Ответ:} $N$-й частичной суммой ряда называется сумма его первых $N$ членов: $S_N = \sum_{n=1}^N a_n$. Суммой ряда $S$ называется предел последовательности частичных сумм при $N \to \infty$: $S = \lim_{N \to \infty} S_N$. Фраза «ряд сходится» означает, что указанный предел существует и является конечным числом.

---

\textbf{77. !!! Сформулируйте необходимый признак сходимости числового ряда. Верно ли об- ратное утверждение? Приведите контрпример.}

\textit{Ответ:} Необходимый признак сходимости: если ряд $\sum a_n$ сходится, то предел его общего члена равен нулю: $\lim_{n \to \infty} a_n = 0$. Обратное утверждение неверно. Контрпример: гармонический ряд $\sum_{n=1}^{\infty} \frac{1}{n}$ расходится, несмотря на то что $\lim_{n \to \infty} \frac{1}{n} = 0$.

---

\textbf{78. !!! Вычислить сумму простейшего эталонного геометрического ряда $\sum_{n=0}^{\infty}q^{n}$ при условии $|q|<1$.}

\textit{Ответ:} Сумма ряда равна:
$$ \sum_{n=0}^{\infty} q^n = \frac{1}{1-q} $$

---

\textbf{79. !!! Сформулируйте критерий Коши для сходимости числового ряда на языке $\epsilon-N$.}

\textit{Ответ:} Числовой ряд $\sum a_n$ сходится тогда и только тогда, когда:
$$ \forall \epsilon > 0 \quad \exists N \in \mathbb{N} \quad \forall n > N \quad \forall p \in \mathbb{N} \implies \left| \sum_{k=n+1}^{n+p} a_k \right| < \epsilon $$

---

\textbf{80. !!! Дайте определение К-го остатка (хвоста) числового ряда $R_{K}$. Как сходимость остатка связана со сходимостью исходного ряда?}

\textit{Ответ:} $K$-м остатком числового ряда называется ряд, полученный отбрасыванием первых $K$ членов: $R_K = \sum_{n=K+1}^{\infty} a_n$. Сходимость остатка прямо эквивалентна сходимости исходного ряда; при этом для сходящегося ряда выполнено $\lim_{K \to \infty} R_K = 0$.

---

\textbf{81. Запишите отрицание критерия Коши (условие расходимости числового ряда) в кван- торах.}

\textit{Ответ:} Условие расходимости ряда в кванторах:
$$ \exists \epsilon > 0 \quad \forall N \in \mathbb{N} \quad \exists n > N \quad \exists p \in \mathbb{N} : \left| \sum_{k=n+1}^{n+p} a_k \right| \ge \epsilon $$

---

\textbf{82. Сформулируйте свойство линейности для сходящихся числовых рядов (сумма рядов и умножение на константу).}

\textit{Ответ:} Если ряды $\sum a_n$ и $\sum b_n$ сходятся к суммам $A$ и $B$ соответственно, то для любых чисел $\alpha, \beta \in \mathbb{R}$ ряд $\sum (\alpha a_n + \beta b_n)$ также сходится, и его сумма равна $\alpha A + \beta B$.

---

\textbf{83. Если ряд $\sum a_{n}$ сходится, а ряд $\sum b_{n}$ расходится, что можно сказать о сходимости ряда $\sum(a_{n}+b_{n})$?}

\textit{Ответ:} Ряд $\sum(a_n + b_n)$ гарантированно расходится.

---

\textbf{84. Изменится ли факт сходимости числового ряда, если мы отбросим первые 1000000 его членов? Изменится ли при этом его сумма?}

\textit{Ответ:} Факт сходимости ряда от этого не изменится. Однако числовое значение суммы ряда изменится (на величину суммы этих отброшенных первых 1 000 000 членов).

---

\textbf{85. Сформулируйте критерий сходимости ряда с неотрицательными членами}

\textit{Ответ:} Ряд с неотрицательными членами $\sum a_n$ ($a_n \ge 0$) сходится тогда и только тогда, когда последовательность его частичных сумм $\{S_N\}$ ограничена сверху.

---

\textbf{86. !!! Сформулируйте признак сравнения числовых рядов в форме неравенств.}

\textit{Ответ:} Пусть для двух рядов с неотрицательными членами выполнено неравенство $0 \le a_n \le b_n$ для всех номеров $n$, начиная с некоторого $n_0$. Тогда:
1) Если ряд $\sum b_n$ сходится, то ряд $\sum a_n$ также сходится. \\
2) Если ряд $\sum a_n$ расходится, то ряд $\sum b_n$ также расходится.

---

\textbf{87. !!! Сформулируйте признак сравнения в предельной форме (через эквивалентность бесконечно малых общего члена при $n\rightarrow\infty$).}

\textit{Ответ:} Если члены рядов $\sum a_n$ и $\sum b_n$ положительны и их общие члены эквивалентны при $n \to \infty$ ($a_n \sim b_n$, то есть $\lim_{n \to \infty} \frac{a_n}{b_n} = 1$), то оба ряда сходятся или расходятся одновременно.

---

\textbf{88. !!! Сформулируйте признак Д'Аламбера сходимости числового ряда в предельной форме.}

\textit{Ответ:} Пусть $a_n > 0$ и существует предел $D = \lim_{n \to \infty} \frac{a_{n+1}}{a_n}$. Тогда:
1) Если $D < 1$, ряд сходится. \\
2) Если $D > 1$ (включая $D = +\infty$), ряд расходится. \\
3) Если $D = 1$, признак не дает ответа.

---

\textbf{89. !!! Запишите формулу общего члена эталонного обобщенного гармонического ряда Дирихле и укажите условия его сходимости.}

\textit{Ответ:} Формула общего члена: $a_n = \frac{1}{n^p}$. Ряд $\sum_{n=1}^{\infty} \frac{1}{n^p}$ сходится при $p > 1$ и расходится при $p \le 1$.

---

\textbf{90. !!! Сформулируйте радикальный признак Коши в предельной форме.}

\textit{Ответ:} Пусть $a_n \ge 0$ и существует предел $C = \lim_{n \to \infty} \sqrt[n]{a_n}$. Тогда:
1) Если $C < 1$, ряд сходится. \\
2) Если $C > 1$, ряд расходится. \\
3) Если $C = 1$, признак не дает ответа.

---

\textbf{91. Какой из двух признаков Д'Аламбера или Коши является более чувствительным («сильным»)? Сформулируйте соответствующее математическое утверждение.}

\textit{Ответ:} Радикальный признак Коши является более сильным (чувствительным), чем признак Д'Аламбера. \\
Математическое утверждение: для любой последовательности положительных чисел $a_n > 0$ справедливо неравенство:
$$ \liminf_{n \to \infty} \frac{a_{n+1}}{a_n} \le \liminf_{n \to \infty} \sqrt[n]{a_n} \le \limsup_{n \to \infty} \sqrt[n]{a_n} \le \limsup_{n \to \infty} \frac{a_{n+1}}{a_n} $$
Отсюда следует, что если предел отношения $\frac{a_{n+1}}{a_n}$ существует и равен $L$, то предел $\sqrt[n]{a_n}$ также существует и равен $L$. Однако предел корня может существовать и давать ответ даже тогда, когда предел отношения не существует.

---

\textbf{92. !!! Дайте определение абсолютной сходимости числового ряда.}

\textit{Ответ:} Числовой ряд $\sum a_n$ называется абсолютно сходящимся, если сходится ряд, составленный из абсолютных величин (модулей) его членов: $\sum |a_n|$.

---

\textbf{93. Сформулируйте свойство абсолютно сходящихся рядов относительно перестановки их членов (теорема Дирихле).}

\textit{Ответ:} Если числовой ряд сходится абсолютно, то любой ряд, полученный из него произвольной перестановкой членов, также сходится абсолютно и имеет ту же самую сумму, что и исходный ряд.

---

\textbf{94. Если ряд $\sum a_{n}$ сходится абсолютно, а ряд $\sum b_{n}$ сходится условно, в какой категории окажется ряд $\sum(a_{n}+b_{n})$?}

\textit{Ответ:} Ряд $\sum(a_n + b_n)$ будет сходиться условно.

---

\textbf{95. !!! Дайте определение условной сходимости числового ряда.}

\textit{Ответ:} Числовой ряд $\sum a_n$ называется условно сходящимся, если сам ряд $\sum a_n$ сходится, но ряд из его модулей $\sum |a_n|$ расходится.

---

\textbf{96. Запишите формулу общего члена классического знакочередующегося гармоническо- го ряда и укажите его точную сумму.}

\textit{Ответ:} Формула общего члена: $a_n = \frac{(-1)^{n-1}}{n}$. Его точная сумма равна:
$$ \sum_{n=1}^{\infty} \frac{(-1)^{n-1}}{n} = \ln 2 $$

---

\textbf{97. Сформулируйте (в кванторах) необходимое условие сходимости для любого знако- переменного числового ряда.}

\textit{Ответ:} Необходимое условие: $\lim_{n \to \infty} a_n = 0$, что в кванторах записывается как:
$$ \forall \epsilon > 0 \quad \exists N \in \mathbb{N} \quad \forall n > N \implies |a_n| < \epsilon $$

---

\textbf{98. !!! Сформулируйте признак Лейбница для знакочередующихся рядов (оба условия, налагаемые на модули членов).}

\textit{Ответ:} Знакочередующийся ряд $\sum_{n=1}^{\infty} (-1)^{n-1} c_n$ (где $c_n > 0$) сходится, если выполнены два условия:
1) Последовательность модулей монотонно не возрастает: $c_{n+1} \le c_n$ для всех $n \in \mathbb{N}$. \\
2) Предел модулей равен нулю: $\lim_{n \to \infty} c_n = 0$.

---

\textbf{99. Сформулируйте следствие из теоремы Лейбница об оценке модуля остатка (погреш- ности) ряда $R_{N}$ после суммирования первых членов.}

\textit{Ответ:} Модуль остатка ряда Лейбница $R_N$ после суммирования первых $N$ членов не превосходит абсолютной величины первого отброшенного члена:
$$ |R_N| \le c_{N+1} $$

---

\textbf{100. Какой знак всегда имеет остаток $R_{N}$ ряда Лейбница? Как он связан со знаком пер- вого отброшенного члена?}

\textit{Ответ:} Знак остатка $R_N$ ряда Лейбница всегда совпадает со знаком первого из отброшенных членов (то есть со знаком $(N+1)$-го члена исходного ряда).

---

\textbf{101. Запишите формулу преобразования Абеля (суммирования по частям) для конечных сумм вида $\sum_{k=1}^{n}a_{k}b_{k}$}

\textit{Ответ:} Обозначим $B_k = \sum_{j=1}^k b_j$ и $B_0 = 0$. Формула имеет вид:
$$ \sum_{k=1}^{n} a_k b_k = a_n B_n + \sum_{k=1}^{n-1} (a_k - a_{k+1}) B_k $$

---

\textbf{102. Сформулируйте лемму Абеля об оценке частичных сумм произведения двух после- довательностей (при условии монотонности одной из них).}

\textit{Ответ:} Если последовательность $\{a_k\}$ неотрицательна и монотонно не возрастает ($a_1 \ge a_2 \ge \dots \ge a_n \ge 0$), а частичные суммы последовательности $\{b_k\}$ ограничены числами $m$ и $M$, то есть $m \le B_k \le M$ (где $B_k = \sum_{j=1}^k b_j$), то справедлива оценка:
$$ a_1 m \le \sum_{k=1}^n a_k b_k \le a_1 M $$

---

\textbf{103. В чем заключается аналогия между преобразованием Абеля в дискретном анализе и формулой интегрирования по частям в непрерывном анализе?}

\textit{Ответ:} Дискретное преобразование Абеля перебрасывает операцию взятия разности (аналог дифференцирования $\Delta a_k = a_{k+1} - a_k$) с одной последовательности на другую путем суммирования (аналог интегрирования $B_k = \sum b_k$). Это в точности соответствует непрерывной формуле $\int u dv = uv - \int v du$.

---

\textbf{104. Сформулируйте признак Дирихле сходимости числового ряда $\sum a_{n}b_{n}$ (условия, на- лагаемые на последовательности $\{a_{n}\}$ и $\{b_{n}\}$ survive).}

\textit{Ответ:} Ряд $\sum a_n b_n$ сходится, если выполнены условия:
1) Последовательность $\{a_n\}$ монотонно стремится к нулю при $n \to \infty$. \\
2) Последовательность частичных сумм ряда $\sum b_n$ ограничена в совокупности: $\exists M > 0 \,\, \forall N \in \mathbb{N} \implies \left| \sum_{n=1}^N b_n \right| \le M$.

---

\textbf{105. Укажите условия (показатель степени р) сходимости и абсолютной сходимости эта- лонного ряда Дирихле $\sum_{n=1}^{\infty}\frac{sin~n}{n^{p}}$}

\textit{Ответ:} Ряд сходится абсолютно при $p > 1$. Ряд сходится условно при $0 < p \le 1$. При $p \le 0$ ряд расходится.

---

\textbf{106. Запишите формулу для вычисления замкнутой суммы косинусов или синусов с воз- растающими аргументами: $\sum_{k=1}^{n}sin~kx$ . Каким свойством обладает эта сумма при $x\ne2\pi m$?}

\textit{Ответ:} Формула замкнутой суммы синусов:
$$ \sum_{k=1}^{n} \sin kx = \frac{\sin \frac{nx}{2} \sin \frac{(n+1)x}{2}}{\sin \frac{x}{2}} $$
При $x \neq 2\pi m$ ($m \in \mathbb{Z}$) знаменатель не обращается в нуль, и эта сумма является ограниченной в совокупности по $n$ величиной: $\left|\sum_{k=1}^n \sin kx \right| \le \frac{1}{|\sin(x/2)|}$.

---

\textbf{107. Сформулируйте признак Абеля сходимости числового ряда $\sum a_{n}b_{n}$}

\textit{Ответ:} Ряд $\sum a_n b_n$ сходится, если выполнены следующие два условия:
1) Числовой ряд $\sum b_n$ сходится. \\
2) Последовательность $\{a_n\}$ является монотонной и ограниченной.

---

\textbf{108. Проведите сравнительный анализ: в чем заключается различие в требованиях к по- ведению монотонного сомножителя $\{a_{n}\}$ в признаках Дирихле и Абеля?}

\textit{Ответ:} В признаке Дирихле от монотонной последовательности $\{a_n\}$ требуется строгое стремление к нулю ($\lim a_n = 0$), в то время как от ряда $\sum b_n$ требуется лишь ограниченность частичных сумм. В признаке Абеля от монотонной последовательности $\{a_n\}$ требуется только ограниченность (стремление к нулю не обязательно), но от ряда $\sum b_n$ требуется полноценная сходимость.

---

\textbf{109. Если ряд $\sum b_{n}$ сходится, а последовательность $\{a_{n}\}$ монотонна и ограничена, обязан ли сходиться ряд из модулей $\sum|a_{n}b_{n}|$?}

\textit{Ответ:} Нет, не обязан. Если ряд $\sum b_n$ сходится условно (например, $b_n = \frac{(-1)^n}{n}$), а последовательность $a_n = 1$ (монотонна и ограничена), то ряд модулей $\sum |a_n b_n| = \sum \frac{1}{n}$ представляет собой расходящийся гармонический ряд.

---

\textbf{110. !!! Запишите определение поточечной сходимости функциональной последователь- ности $\{f_{n}(x)\}$ к предельной функции $f(x)$ на множестве Х (в кванторах).}

\textit{Ответ:} Последовательность $\{f_n(x)\}$ сходится поточечно к $f(x)$ на множестве $X$, если:
$$ \forall x \in X \quad \forall \epsilon > 0 \quad \exists N = N(x, \epsilon) \in \mathbb{N} \quad \forall n > N \implies |f_n(x) - f(x)| < \epsilon $$

---

\textbf{111. !!! Запишите определение равномерной сходимости functional последователь- ности $\{f_{n}(x)\}\rightrightarrows f(x)$ на множестве Х (в кванторах). В чем разница в порядке следования кванторов по сравнению с поточечным случаем?}

\textit{Ответ:} Определение равномерной сходимости:
$$ \forall \epsilon > 0 \quad \exists N = N(\epsilon) \in \mathbb{N} \quad \forall x \in X \quad \forall n > N \implies |f_n(x) - f(x)| < \epsilon $$
Разница в порядке кванторов: в поточечном определении квантор $\forall x$ стоит на первом месте, из-за чего номер $N(x, \epsilon)$ зависит от конкретной точки $x$. В равномерном определении квантор $\forall x$ перенесен вглубь (после $\exists N$), поэтому выбранный номер $N(\epsilon)$ является общим и единым для всех точек множества $X$.

---

\textbf{112. Опишите геометрическую интерпретацию равномерной сходимости через критерий -полосы графика предельной функции.}

\textit{Ответ:} Геометрически равномерная сходимость означает, что начиная с некоторого номера $N$, графики всех последующих функций $y = f_n(x)$ целиком укладываются внутрь полосы шириной $2\epsilon$, симметрично окружающей график предельной функции $y = f(x)$, то есть ограниченной сверху линией $y = f(x) + \epsilon$ и снизу линией $y = f(x) - \epsilon$.

---

\textbf{113. Сформулируйте супремум-критерий равномерной сходимости функциональной по- следовательности через величину $\rho_{n}=sup_{x\in X}|f_{n}(x)-f(x)|$.}

\textit{Ответ:} Последовательность $\{f_n(x)\}$ сходится равномерно к $f(x)$ на множестве $X$ тогда и только тогда, когда числовая последовательность супремумов ее отклонений стремится к нулю:
$$ \lim_{n \to \infty} \rho_n = 0, \quad \text{где} \quad \rho_n = \sup_{x \in X} |f_n(x) - f(x)| $$

---

\textbf{114. Справедливо ли утверждение: «Если последовательность непрерывных функций схо- дится поточечно к непрерывной функции, то эта сходимость обязательно равномер- на»? Приведите краткий комментарий.}

\textit{Ответ:} Нет, утверждение неверно. Например, последовательность $f_n(x) = x^n$ на полуинтервале $[0, 1)$ поточечно сходится к непрерывному нулю $f(x) \equiv 0$, однако сходимость неравномерная, так как $\sup_{x \in [0, 1)} |x^n - 0| = 1 \not\to 0$. Для равномерности требуются дополнительные условия (например, компактность множества и монотонность сходимости по теореме Дини).

---

\textbf{115. !!! Сформулируйте критерий Коши равномерной сходимости функциональной по- следовательности на множестве Х (на языке $\epsilon-\Lambda$).}

\textit{Ответ:} Последовательность $\{f_n(x)\}$ сходится равномерно на множестве $X$ тогда и только тогда, когда:
$$ \forall \epsilon > 0 \quad \exists N \in \mathbb{N} \quad \forall n, m > N \quad \forall x \in X \implies |f_n(x) - f_m(x)| < \epsilon $$

---

\textbf{116. !!! Сформулируйте критерий Коши равномерной сходимости функционального ряда $\sum_{n=1}^{\infty}u_{n}(x)$ на множестве Х.}

\textit{Ответ:} Функциональный ряд $\sum u_n(x)$ сходится равномерно на множестве $X$ тогда и только тогда, когда:
$$ \forall \epsilon > 0 \quad \exists N \in \mathbb{N} \quad \forall n > N \quad \forall p \in \mathbb{N} \quad \forall x \in X \implies \left| \sum_{k=n+1}^{n+p} u_k(x) \right| < \epsilon $$

---

\textbf{117. В чем заключается отличие критерия Коши от супремум-критерия с точки зрения практического применения?}

\textit{Ответ:} Критерий Коши позволяет установить факт равномерной сходимости функциональной последовательности или ряда «внутренним» образом, без предварительного нахождения явного аналитического вида предельной функции (или суммы ряда), в то время как супремум-критерий требует обязательного знания предельной функции $f(x)$.

---

\textbf{118. Дайте определение равномерной сходимости функционального ряда $\sum_{n=1}^{\infty}u_{n}(x)$ на множестве Х через поведение последовательности его остатков $R_{n}(x)$.**}

\textit{Ответ:} Функциональный ряд сходится равномерно на множестве $X$, если последовательность его остаточных функций $R_n(x) = \sum_{k=n+1}^{\infty} u_k(x)$ равномерно стремится к нулю на этом множестве при $n \to \infty$, то есть:
$$ \lim_{n \to \infty} \left( \sup_{x \in X} |R_n(x)| \right) = 0 $$

---

\textbf{119. !!! Сформулируйте мажорантный признак Вейерштрасса равномерной сходимости функционального ряда.}

\textit{Ответ:} Если члены функционального ряда $\sum u_n(x)$ на множестве $X$ ограничены по модулю членами сходящегося положительного числового ряда $\sum a_n$, то есть $\forall x \in X$ и $\forall n \in \mathbb{N}$ выполнено $|u_n(x)| \le a_n$, то данный функциональный ряд сходится на множестве $X$ абсолютно и равномерно.

---

\textbf{120. Справедливо ли утверждение: «Если функциональный ряд сходится равномерно на множестве Х, то он сходится на этом множестве и абсолютно»? Приведите контр- пример.}

\textit{Ответ:} Нет, не справедливо. Контрпример: ряд $\sum_{n=1}^{\infty} \frac{(-1)^{n-1}}{n+x^2}$ на всей оси $\mathbb{R}$. По признаку Лейбница он сходится равномерно, однако ряд из его модулей $\sum_{n=1}^{\infty} \frac{1}{n+x^2}$ расходится в каждой точке по признаку сравнения с гармоническим рядом.

---

\textbf{121. !!! Сформулируйте теорему о непрерывности суммы функционального ряда, члены которого непрерывны на множестве Х.}

\textit{Ответ:} Если все члены функционального ряда $\sum u_n(x)$ непрерывны на множестве $X$ и ряд сходится к функции $S(x)$ равномерно на этом множестве, то его сумма $S(x)$ также является непрерывной функцией на множестве $X$.

---

\textbf{122. Запишите формулировку теоремы о возможности перестановки знака предела по переменной и знака бесконечной суммы ряда: $lim_{x\rightarrow x_{0}}\sum u_{n}(x)$.}

\textit{Ответ:} Если функциональный ряд $\sum u_n(x)$ сходится равномерно в окрестности точки $x_0$ и для каждого члена существует конечный предел $\lim_{x \to x_0} u_n(x) = c_n$, то числовой ряд $\sum c_n$ сходится, и справедливо равенство:
$$ \lim_{x \to x_0} \sum_{n=1}^{\infty} u_n(x) = \sum_{n=1}^{\infty} \lim_{x \to x_0} u_n(x) = \sum_{n=1}^{\infty} c_n $$

---

\textbf{123. Приведите пример функциональной последовательности непрерывных функций, пре- дел которой является разрывной функцией}

\textit{Ответ:} Последовательность $f_n(x) = x^n$ на замкнутом сегменте $[0, 1]$. Каждая функция $f_n(x)$ непрерывна, однако предельная функция имеет разрыв первого рода в точке $x = 1$:
$$ f(x) = \lim_{n \to \infty} x^n = \begin{cases} 0, & 0 \le x < 1 \\ 1, & x = 1 \end{cases} $$

---

\textbf{124. !!! Сформулируйте теорему о почленном интегрировании функционального ряда. Какие требования предъявляются к промежутку интегрирования?}

\textit{Ответ:} Если члены функционального ряда $\sum u_n(x)$ непрерывны на сегменте $[a, b]$ и ряд сходится на нем равномерно к функции $S(x)$, то его можно почленно интегрировать:
$$ \int_a^b S(x) dx = \sum_{n=1}^{\infty} \int_a^b u_n(x) dx $$
Требование к промежутку интегрирования: он должен быть конечным замкнутым отрезком (сегментом).

---

\textbf{125. !!! Сформулируйте теорему о почленном дифференцировании функционального ря- да. Сходимость какого именно ряда (исходного или ряда производных) обязана быть равномерной?}

\textit{Ответ:} Пусть функции $u_n(x)$ непрерывно дифференцируемы на промежутке $X$. Если сам ряд $\sum u_n(x)$ сходится хотя бы в одной точке $x_0 \in X$, а ряд его производных $\sum u_n'(x)$ сходится равномерно на $X$, то исходный ряд сходится равномерно на всем промежутке $X$, его сумма $S(x)$ дифференцируема, и выполняется равенство $S'(x) = \sum u_n'(x)$. Равномерной обязана быть сходимость \textbf{ряда производных}.

---

\textbf{126. Запишите уравнение почленного дифференцирования ряда: $\frac{d}{dx}\sum_{n=1}^{\infty}u_{n}(x)=...$}

\textit{Ответ:}
$$ \frac{d}{dx}\sum_{n=1}^{\infty} u_n(x) = \sum_{n=1}^{\infty} \frac{d}{dx} u_n(x) = \sum_{n=1}^{\infty} u_n'(x) $$

---

\textbf{127. !!! Дайте определение степенного ряда по степеням $(x-x_{0})$.}

\textit{Ответ:} Степенным рядом по степеням $(x-x_0)$ называется функциональный ряд вида:
$$ \sum_{n=0}^{\infty} a_n (x - x_0)^n = a_0 + a_1(x-x_0) + a_2(x-x_0)^2 + \dots $$
где $a_n \in \mathbb{R}$ — числовые коэффициенты ряда, а $x_0$ — фиксированный центр ряда.

---

\textbf{128. !!! Сформулируйте первую теорему Абеля о структуре области сходимости степен- ного ряда.}

\textit{Ответ:} Если степенной ряд $\sum a_n x^n$ сходится в некоторой точке $x_1 \neq 0$, то он сходится абсолютно для всех точек $x$, удовлетворяющих условию $|x| < |x_1|$. Если ряд расходится в точке $x_2$, то он расходится во всех точках $x$, для которых $|x| > |x_2|$.

---

\textbf{129. !!! Какую геометрическую форму имеет область сходимости степенного ряда на ве- щественной прямой R?}

\textit{Ответ:} На вещественной прямой область сходимости всегда представляет собой симметричный интервал с центром в точке $x_0$, то есть $(x_0 - R, x_0 + R)$ (интервал сходимости), к которому в зависимости от конкретного ряда могут добавляться или не добавляться граничные точки $x_0 \pm R$. В вырожденных случаях область сужается до единственной точки $\{x_0\}$ или совпадает со всей прямой $\mathbb{R}$.

---

\textbf{130. !!! Запишите формулу Даламбера для вычисления радиуса сходимости степенного ряда через отношение коэффициентов. При каких условиях она применима?}

\textit{Ответ:} Формула Даламбера:
$$ R = \lim_{n \to \infty} \left| \frac{a_n}{a_{n+1}} \right| $$
Она применима при условии, что данный предел (конечный или бесконечный) существует, и коэффициенты ряда $a_n$, начиная с некоторого номера, отличны от нуля.

---

\textbf{131. !!! Сформулируйте теорему Коши-Адамара. Запишите формулу радиуса сходимости.}

\textit{Ответ:} Для любого степенного ряда существует радиус сходимости $R$, который определяется через верхний предел последовательности коэффициентов по формуле Коши-Адамара:
$$ \frac{1}{R} = \limsup_{n \to \infty} \sqrt[n]{|a_n|} $$
При этом принимают $R = 0$, если верхний предел равен $+\infty$, и $R = +\infty$, если предел равен $0$.

---

\textbf{132. !!! Какой диагноз ставится степенному ряду, если по формуле Коши-Адамара полу- чено: а) $R=0$, б) $R=+\infty$?}

\textit{Ответ:}
а) Если $R=0$, ряд абсолютно расходится во всей области, кроме единственной точки $x = x_0$ (в центре сходимости сумма равна $a_0$). \\
б) Если $R=+\infty$, ряд сходится абсолютно и равномерно на любом конечном замкнутом отрезке вещественной прямой $\mathbb{R}$.

---

\textbf{133. !!! Сформулируйте теорему о характере сходимости степенного ряда на любом за- мкнутом отрезке [а, 3], строго вложенном внутрь интервала сходимости (-R, R).}

\textit{Ответ:} На любом замкнутом отрезка $[\alpha, \beta]$, который строго лежит внутри интервала сходимости $(x_0 - R, x_0 + R)$, степенной ряд сходится абсолютно и равномерно.

---

\textbf{134. !!! Запишите теоремы о непрерывности суммы, почленном интегрировании и почлен- ном дифференцировании степенного ряда внутри его интервала сходимости.}

\textit{Ответ:} Внутри интервала сходимости $(x_0 - R, x_0 + R)$:
1) \textbf{Непрерывность:} Функция суммы степенного ряда $S(x)$ является непрерывной. \\
2) \textbf{Интегрирование:} Степенной ряд можно почленно интегрировать по любому отрезку, лежащему внутри интервала сходимости. \\
3) \textbf{Дифференцирование:} Степенной ряд можно почленно дифференцировать локально любое число раз внутри интервала сходимости.

---

\textbf{135. !!! Меняется ли величина радиуса сходимости при формальном почленном диф- ференцировании или интегрировании степенного ряда?}

\textit{Ответ:} Нет, величина радиуса сходимости $R$ при почленном дифференцировании или интегрировании степенного ряда не изменяется.

---

\textbf{136. Дайте определение аналитичности вещественной функции $f(x)$ в точке $x_{0}$ (по Т. Tao).}

\textit{Ответ:} Вещественная функция $f(x)$ называется аналитической в точке $x_0$, если в некоторой открытой окрестности $(x_0 - r, x_0 + r)$ этой точки функция может быть представлена в виде сходящегося к ней степенного ряда по степеням $(x-x_0)$.

---

\textbf{137. Укажите теоретико-множественную иерархию между классами функций $C^{0}$, $C^{1}$ $C^{\infty}$ и $C^{\omega}$ на интервале (а, ). Какой класс вложен глубже всех?}

\textit{Ответ:} Иерархия включения классов имеет вид:
$$ C^\omega(a,b) \subset C^\infty(a,b) \subset \dots \subset C^1(a,b) \subset C^0(a,b) $$
Глубже всех вложен класс аналитических функций $C^\omega(a,b)$.

---

\textbf{138. Сформулируйте необходимое условие аналитичности функции в точке.}

\textit{Ответ:} Необходимым условием аналитичности функции в точке является ее бесконечная дифференцируемость в этой точке (функция должна принадлежать классу $C^\infty$).

---

\textbf{139. !!! Запишите формулу для вычисления коэффициентов Тейлора $a_{n}$ сходящегося сте- пенного ряда через значения производных функции в центре ряда.}

\textit{Ответ:} Коэффициенты вычисляются по формуле:
$$ a_n = \frac{f^{(n)}(x_0)}{n!} $$

---

\textbf{140. !!! Напишите формальный вид ряда Тейлора для функции $f(x)$ в окрестности точки $x_{0}$.}

\textit{Ответ:} Формальный вид ряда Тейлора:
$$ f(x) \sim \sum_{n=0}^{\infty} \frac{f^{(n)}(x_0)}{n!} (x - x_0)^n $$

---

\textbf{141. В каком случае ряд Тейлора называется рядом Маклорена? Укажите центр разло- жения.}

\textit{Ответ:} Ряд Тейлора называется рядом Маклорена, если разложение функции строится в окрестности точки $x_0 = 0$ (центр разложения находится в начале координат).

---

\textbf{142. Запишите выражение для остаточного члена ряда Тейлора $R_{n}(x)$ в форме Лагранжа и в форме Коши.}

\textit{Ответ:} Пусть $\xi = x_0 + \theta(x - x_0)$ при $\theta \in (0, 1)$.
$$ \text{Форма Лагранжа:} \quad R_n(x) = \frac{f^{(n+1)}(\xi)}{(n+1)!} (x - x_0)^{n+1} $$
$$ \text{Форма Коши:} \quad R_n(x) = \frac{f^{(n+1)}(\xi)}{n!} (1 - \theta)^n (x - x_0)^{n+1} $$

---

\textbf{143. Сформулируйте критерий сходимости ряда Тейлора к самой функции через предел последовательности остаточных членов при $n\rightarrow\infty$}

\textit{Ответ:} Ряд Тейлора бесконечно дифференцируемой функции $f(x)$ сходится к значению функции в данной точке $x$ тогда и только тогда, когда предел последовательности остаточных членов равен нулю при $n \to \infty$:
$$ \lim_{n \to \infty} R_n(x) = 0 $$

---

\textbf{144. Запишите достаточное условие аналитичности функции на отрезке через ограниче- ние роста её производных.}

\textit{Ответ:} Если все производные бесконечно дифференцируемой функции $f(x)$ ограничены в совокупности на отрезке $[a, b]$ одной и той же константой $M$ (то есть $\forall n \in \mathbb{N}_0 \,\, \forall x \in [a, b] \implies |f^{(n)}(x)| \le M$), то функция аналитична на этом отрезке.

---

\textbf{145. !!! Запишите четыре аксиомы скалярного произведения в предгильбертовом про- странстве вещественных функций.}

\textit{Ответ:} Для любых элементов $f, g, h$ и чисел $\alpha, \beta \in \mathbb{R}$ скалярное произведение $\langle f, g \rangle$ должно удовлетворять аксиомам:
1) \textbf{Линейность:} $\langle \alpha f + \beta g, h \rangle = \alpha \langle f, h \rangle + \beta \langle g, h \rangle$ \\
2) \textbf{Симметричность:} $\langle f, g \rangle = \langle g, f \rangle$ \\
3) \textbf{Неотрицательность:} $\langle f, f \rangle \ge 0$ \\
4) \textbf{Строгая определенность:} $\langle f, f \rangle = 0 \iff f \equiv 0$ (равенство нулю тождественно).

---

\textbf{146. !!! Сформулируйте неравенство Коши-Буняковского для пространства $L_{2}[a,b]$ в ин- тегральной форме.}

\textit{Ответ:} Интегральная форма неравенства:
$$ \left( \int_a^b f(x)g(x) dx \right)^2 \le \left( \int_a^b f^2(x) dx \right) \left( \int_a^b g^2(x) dx \right) $$

---

\textbf{147. В чем заключается проблема строгости аксиомы положительной определенности $(\langle f,f\rangle=0\Rightarrow f\equiv0)$ для интеграла Римана? Как это испраить введением фактор- пространства?}

\textit{Ответ:} Проблема заключается в том, что интеграл Римана от квадрата функции $\int_a^b f^2(x)dx$ может быть равен нулю для функций, не равных нулю тождественно (например, функция равна нулю везде, кроме конечного числа изолированных точек). Это исправляется переходом к фактор-пространству по подпространству функций, эквивалентных нулю (равных нулю почти всюду), где элементами пространства становятся классы эквивалентных функций.

---

\textbf{148. !!! Дайте определения ортогональной и ортонормированной систем функций в про- странстве $L_{2}[a,b]$.}

\textit{Ответ:} Система функций $\{\phi_n(x)\}$ называется ортогональной, если скалярное произведение любых двух различных элементов равно нулю: $\langle \phi_n, \phi_m \rangle = 0$ при $n \neq m$. Система называется ортонормированной, если она ортогональна и норма каждой функции равна единице, то есть:
$$ \langle \phi_n, \phi_m \rangle = \delta_{nm} = \begin{cases} 1, & n = m \\ 0, & n \neq m \end{cases} $$

---

\textbf{149. !!! Запишите общую формулу вычисления коэффициентов Фурье $c_{n}$ для функции $f(x)$ по произвольной ортогональной системе $\{\phi_{n}(x)\}$.}

\textit{Ответ:} Коэффициенты вычисляются по формуле:
$$ c_n = \frac{\langle f, \phi_n \rangle}{\langle \phi_n, \phi_n \rangle} = \frac{\int_a^b f(x)\phi_n(x)dx}{\int_a^b \phi_n^2(x)dx} $$

---

\textbf{150. Сформулируйте геометрический смысл неравенства Бесселя и равенства Парсеваля в терминах бесконечномерной теоремы Пифагора.}

\textit{Ответ:} Неравенство Бесселя означает, что сумма квадратов длин проекций вектора на взаимно перпендикулярные оси подпространства не превышает квадрата полной длины самого вектора. Равенство Парсеваля означает, что для базиса сумма квадратов координат вектора в точности равна квадрату его длины, что является обобщением теоремы Пифагора на бесконечномерный случай.

---

\textbf{151. Запишите интегральное представление №-й частичной суммы $S_{N}(x)$ тригонометри- ческого ряда Фурье через свертку с ядром Дирихле.}

\textit{Ответ:} Интегральное представление (интеграл Дирихле):
$$ S_N(x) = \frac{1}{\pi} \int_{-\pi}^{\pi} f(x+u)D_N(u)du = \frac{1}{\pi} \int_{-\pi}^{\pi} f(t)D_N(t-x)dt $$

---

\textbf{152. Выпишите компактную аналитическую формулу ядра Дирихле $D_{N}(u)$ через отно- шение синусов и перечислите три его свойства (нормировка, четность, осцилляция).}

\textit{Ответ:} Аналитическая формула ядра Дирихле:
$$ D_N(u) = \frac{\sin\left((N+\frac{1}{2})u\right)}{2\sin\frac{u}{2}} $$
Свойства:
1) \textbf{Нормировка:} $\frac{1}{\pi}\int_{-\pi}^{\pi} D_N(u) du = 1$. \\
2) \textbf{Четность:} $D_N(-u) = D_N(u)$. \\
3) \textbf{Осцилляция:} ядро интенсивно колеблется с высокой частотой при больших $N$, концентрируя основную массу площади около точки $u=0$.

---

\textbf{153. Сформулируйте лемму Римана-Лебега об асимптотическом поведении осциллиру- ющих интегралов при частоте $\lambda\rightarrow+\infty$}

\textit{Ответ:} Если функция $f(x)$ интегрируема на сегменте $[a, b]$, то при неограниченном росте частоты осцилляций интеграл стремится к нулю:
$$ \lim_{\lambda \to +\infty} \int_a^b f(x) \sin(\lambda x) dx = 0 \quad \text{и} \quad \lim_{\lambda \to +\infty} \int_a^b f(x) \cos(\lambda x) dx = 0 $$

---

\textbf{154. Сформулируйте условие Липшица порядка а $(0<\alpha\le1)$ для функции в точке $x_{0}.$ Какой порядок гарантирует обычная дифференцируемость?}

\textit{Ответ:} Функция удовлетворяет условию Липшица порядка $\alpha$ в точке $x_0$, если в некоторой ее окрестности выполняется неравенство $|f(x) - f(x_0)| \le L |x - x_0|^\alpha$, где $L > 0$. Обычная дифференцируемость функции в точке гарантирует выполнение условия Липшица первого порядка ($\alpha = 1$).

---

\textbf{155. !!! Запишите формулировку признака Дини поточечной сходимости тригонометри- ческого ряда Фурье.}

\textit{Ответ:} Если для функции $f(x)$ в некоторой точке $x$ сходится интеграл Дини:
$$ \int_0^\delta \frac{|f(x+u) + f(x-u) - 2f(x)|}{u} du < +\infty $$
(для некоторого $\delta > 0$), то тригонометрический ряд Фурье функции $f$ в этой точке сходится локально к значению $f(x)$.

---

\textbf{156. Сходится ли ряд Фурье непрерывной функции в каждой её точке?}

\textit{Ответ:} Нет, вообще говоря, не обязан. Существует классический контрпример Дюбуа-Реймона непрерывной функции, ряд Фурье которой расходится в фиксированной точке. Для поточечной сходимости непрерывной функции требуется дополнительная гладкость (например, кусочная дифференцируемость или выполнение условия Липшица/Дини).

---

\textbf{157. !!! Сформулируйте теорему о поведении и значении суммы ряда Фурье в точках разрыва первого рода.}

\textit{Ответ:} Если функция $f(x)$ кусочно-гладкая, то в точке $x_0$ ее разрыва первого рода тригонометрический ряд Фурье сходится к среднему арифметическому значению ее односторонних пределов:
$$ S(x_0) = \frac{f(x_0 - 0) + f(x_0 + 0)}{2} $$

---
---

\textbf{158. Запишите формулу числового ряда Лейбница для числа $\pi$, получаемую при разложении в ряд Фурье периодической прямоугольной волны.}

\textit{Ответ:} Формула ряда Лейбница:
$$ \frac{\pi}{4} = 1 - \frac{1}{3} + \frac{1}{5} - \frac{1}{7} + \dots = \sum_{n=0}^{\infty} \frac{(-1)^n}{2n+1} $$

---

\textbf{159. !!! Запишите аксиому косой симметрии (эрмитовости) для комплексного скалярного произведения в пространстве $L_{2}^{\mathbb{C}}[a,b]$.}

\textit{Ответ:} Аксиома косой симметрии (комплексного сопряжения):
$$ \langle f, g \rangle = \overline{\langle g, f \rangle} $$

---

\textbf{160. !!! Запишите формулу вычисления комплексного коэффициента Фурье $c_{n}$ для функции $f(x)$ на сегменте $[-\pi,\pi]$ по системе экспонент.}

\textit{Ответ:} Комплексный коэффициент Фурье вычисляется по формуле:
$$ c_n = \frac{1}{2\pi} \int_{-\pi}^{\pi} f(x) e^{-inx} dx $$

---

\textbf{161. Сформулируйте теорему Планшереля (равенство Парсеваля) сходимости ряда Фурье в терминах бесконечномерной теоремы Пифагора.}

\textit{Ответ:} Если функция $f(x) \in L_2[-\pi, \pi]$, то для её коэффициентов Фурье выполняется равенство Парсеваля:
$$ \frac{1}{\pi} \int_{-\pi}^{\pi} |f(x)|^2 dx = \frac{a_0^2}{2} + \sum_{n=1}^{\infty} (a_n^2 + b_n^2) $$
Для комплексной формы:
$$ \frac{1}{2\pi} \int_{-\pi}^{\pi} |f(x)|^2 dx = \sum_{n=-\infty}^{\infty} |c_n|^2 $$

\end{document}